\subsection{Method}
\label{sec:method}

The main method, the \textbf{Failure-Trace Allocator (FTA)}, is an
answer-selection layer applied \emph{after} candidate generation
(Figure~\ref{fig:method_overview}).
It does not generate new candidates and does not issue additional inference
calls: it selects a final answer from already-produced frontier and external
outputs under the same per-method budget cap.
Its two gates---the \emph{Low-Depth Guard} and the
\emph{External-Consensus Fallback}---fire in strict priority order; at most
one gate acts per example.
All runtime decisions are gold-free: no gold answer, exact-match signal, or
correctness label is accessed.
Implementation-name mapping is retained in Table~\ref{tab:name_mapping}.

At each example, FTA consumes three classes of runtime-legal inputs:
the frontier answer candidate, the three evaluated external-baseline answers,
and failure-trace metadata emitted during normal frontier execution.
The output is a single final answer plus an explicit policy outcome
(\texttt{fix2}, \texttt{fix4}, or frontier retained), so every decision can be
audited after the run using only logged runtime fields.
This interface is intentionally narrow: the main method does not depend on
learned scores or post-hoc labels.

From an intelligent-systems perspective, FTA can be read as an
instance-level selective deferral policy over a fixed candidate pool:
rather than learning a new generator, it uses runtime-visible
failure-trace and agreement signals to decide whether to retain the
frontier answer or defer to an external-consensus answer under a
matched-budget inference contract.
This places the method in the same conceptual space as selective
prediction, dynamic ensemble selection, and cost-aware LLM routing, while
keeping the decision rule deterministic and auditable.

\paragraph{Frontier generator and runtime metadata.}
The frontier method (Section~\ref{sec:terminology}) produces each candidate by exploring
multiple reasoning \emph{branches}, where a branch is one candidate trajectory and its
\emph{depth} is the number of expansion steps from the root.
As a side-effect of normal inference, the frontier logs four runtime-visible metadata
fields that FTA consumes.
\texttt{frontier\_support} is the integer count of branches that converge on the selected
answer; a value of zero or one indicates a low-support, single-branch selection.
\texttt{candidate\_pool\_answer\_group\_count} is the integer count of distinct answer
equivalence classes across all branches explored, with higher values indicating greater
internal disagreement.
\texttt{override\_reason} is a string set by the frontier's internal selection step;
the values \texttt{"single\_weak\_frontier\_branch"} (at most one branch supports the
selected answer) and \texttt{"direct\_frontier\_agree"} (the frontier answer matches the
frontier's own internal reference selection, indicating internal self-consistency) are the
two states relevant to the LDG and ECO gates respectively.
\texttt{direct\_frontier\_agree} is a boolean indicator of the same internal-consistency
condition.
All fields are recorded during ordinary inference without gold-label access; FTA uses them
exactly as logged, making the decision boundary auditable from the run records.

Conceptually, FTA uses three abstract predicates.
\textsc{WeakFrontier}$(x)$ is true when the frontier trace indicates low support
or a single-branch selection.
\textsc{ExternalMajority}$(x)$ returns an external answer when at least two of
L1, S1, and TALE agree, with a deterministic fallback only when all three
differ.
\textsc{ExternalConsensusDisagree}$(x)$ is true when L1, S1, and TALE
unanimously agree with one another and disagree with the frontier answer.
The exact field-level realizations of these predicates are given in
Table~\ref{alg:fix24} and Table~\ref{tab:policy_rules}; the raw log-field
names are implementation details documented in the reproducibility appendix.

The \emph{Low-Depth Guard} is motivated by a structural reliability pattern.
When the frontier search terminates with a single weak branch---signalled by
logged branch-trace metadata---the resulting candidate is treated as evidence
of under-exploration.
In such cases, routing to an external majority answer provides an independent
cross-policy signal without spending extra compute.
The \emph{External-Consensus Fallback} addresses a complementary pattern: if
the frontier's selection is internally self-consistent yet all three external
methods unanimously prefer a different answer, that unanimous disagreement is
treated as high-confidence evidence that the frontier selection may be wrong.
Partial external agreement (one or two externals disagreeing) is intentionally
not used in the main method because it did not provide validation-grade evidence
in the evaluation setting.

\paragraph{Selective-recovery proposition.}
Let $S$ be a runtime-identifiable subset of examples.
Let $e_F(S)$ be frontier error rate on $S$ and $e_M(S)$ be external-majority
error rate on $S$.
If $e_F(S) > e_M(S)$, then replacing frontier with external majority on $S$
improves expected accuracy by
\[
  P(S)\,\bigl(e_F(S)-e_M(S)\bigr).
\]
The Low-Depth Guard (LDG) and External-Consensus Fallback (ECO) are
gold-free empirical heuristics intended to identify such subsets without using
correctness labels at runtime.

The control flow follows a strict precedence chain
(Figure~\ref{fig:method_overview}, Table~\ref{tab:policy_rules}).
The Low-Depth Guard is checked first; when it fires, the policy routes to an
external majority answer and stops.
This priority reflects the interpretation of the weak-frontier trace as a
generation-quality warning: when exploration appears narrow, the policy does
not additionally check external consensus.
If LDG does not fire, the External-Consensus Fallback is checked; when that
condition fires, the policy routes to the external unanimity answer.
If neither gate fires, the frontier answer is retained.
The default therefore preserves the original frontier method on all examples
where neither high-confidence failure trace is present.
Conflict handling is deterministic: at most one gate can act, and external
majority ties are resolved by the fixed priority order documented in
Table~\ref{tab:policy_rules}.
That fallback is retained for reproducibility rather than claimed as optimal;
the sensitivity noted in Table~\ref{tab:param_sensitivity} shows that
alternative tie orders change Aggregate-720 replay by less than one
percentage point.

All trigger features are runtime-legal and gold-free by construction.
They are computed only from model outputs and frontier trace metadata that are
already visible at inference time, never from gold answers, exact-match
signals, example identifiers, or artifact paths.
Table~\ref{tab:feature_legality} gives the complete legality inventory,
including allowed inputs and explicitly excluded fields.
This legality boundary is central to the method definition, not only a
reporting constraint.

Relative to L1, S1, TALE, and the frontier baseline, FTA plays a different
role in the stack.
L1, S1, TALE, and frontier each provide candidate answers under the same
matched per-method budget; FTA then decides whether to keep the frontier
answer or override it using external agreement signals.
This separation between candidate generation and final allocation is what
allows the policy to improve selection quality without additional model calls.
Observed gains come from switching the final answer on a small
trace-identified subset of examples rather than from increasing inference
budget.

\subsection{Low-Depth Guard (LDG)}
\label{sec:fix2}

The Low-Depth Guard detects frontier answers arising from structurally weak
search branches and replaces them with an external majority vote.
The trigger relies exclusively on fields logged during the frontier run:
the override-reason indicator and the candidate-pool group count.

\subsection{External-Consensus Fallback (ECO)}
\label{sec:fix4}

The External-Consensus Fallback identifies examples where the frontier agreed
with its own internal reserve selection, yet all three external baselines
unanimously chose a different answer.
When this unanimous disagreement holds, the policy substitutes the external
consensus answer.

\subsection{Combined Policy Summary}
\label{sec:fix24_algorithm}

Table~\ref{alg:fix24} gives the complete pseudocode using exact runtime field
names.
Parameter settings and implementation-name mapping are documented in the
appendix reproducibility tables.

\begin{table}[t]
\centering
\caption{Failure-Trace Allocator (FTA) combined policy.
All inputs are gold-free runtime fields.
Named constants and name mapping are documented in
Table~\ref{tab:name_mapping} and Table~\ref{tab:policy_rules}.}
\label{alg:fix24}
\scriptsize
\begin{tabular}{p{0.06\textwidth} p{0.88\textwidth}}
\toprule
\textbf{Step} & \textbf{Operation} \\
\midrule
\multicolumn{2}{l}{\textit{Inputs per example}} \\
& \texttt{frontier\_answer}, \texttt{result\_metadata} (from frontier run logs) \\
& \texttt{l1\_answer}, \texttt{s1\_answer}, \texttt{tale\_answer} (external baseline outputs) \\
& \texttt{override\_reason}, \texttt{frontier\_support} $\in \mathbb{Z}_{\geq 0}$, \texttt{candidate\_pool\_answer\_group\_count} $\in \mathbb{Z}_{>0}$ \\
\midrule
\multicolumn{2}{l}{\textit{Step 1 --- Low-Depth Guard (LDG)}} \\
1a & \textbf{if} \texttt{override\_reason} contains \texttt{"single\_weak\_frontier\_branch"} \\
   & \quad \textbf{or} (\texttt{frontier\_support} $= 0$ \textbf{and} \texttt{candidate\_pool\_answer\_group\_count} $\geq 2$) \\
   & \textbf{then:} \\
1b & \quad Compute \texttt{ext} $\leftarrow$ majority answer among \{\texttt{l1}, \texttt{s1}, \texttt{tale}\} (at least 2 agree); \\
   & \quad if no majority, use priority order TALE $\succ$ S1 $\succ$ L1 \\
1c & \quad \textbf{if} \texttt{ext} exists \textbf{and} \texttt{ext} $\neq$ \texttt{frontier\_answer}: \\
   & \quad\quad \textbf{return} \texttt{ext}, \texttt{policy} $\leftarrow$ \texttt{fix2} \quad\quad \textit{(guard fires; stop)} \\
\midrule
\multicolumn{2}{l}{\textit{Step 2 --- External-Consensus Fallback (ECO)}} \\
2a & \textbf{if} \texttt{override\_reason} $=$ \texttt{"direct\_frontier\_agree"} \\
   & \quad \textbf{and} \texttt{l1\_answer} $=$ \texttt{s1\_answer} $=$ \texttt{tale\_answer} \\
   & \quad \textbf{and} \texttt{l1\_answer} $\neq$ \texttt{frontier\_answer} \\
   & \textbf{then:} \\
2b & \quad \textbf{return} \texttt{l1\_answer}, \texttt{policy} $\leftarrow$ \texttt{fix4} \quad\quad \textit{(consensus fallback fires; stop)} \\
\midrule
\multicolumn{2}{l}{\textit{Step 3 --- Default}} \\
3  & \textbf{return} \texttt{frontier\_answer}, \texttt{policy} $\leftarrow$ \texttt{original} \\
\midrule
\multicolumn{2}{l}{\textit{Output}} \\
& Selected answer and \texttt{policy} label; no gold labels accessed at any step. \\
\bottomrule
\end{tabular}
\end{table}

The exact field-level trigger conditions are also tabulated in
Table~\ref{tab:policy_rules} for reproducibility.
